\input{../tex-inputs/latex-leseansicht-vorspann}

\section[Gustav Schwarzkopf an Arthur Schnitzler, 16. 6. 1928]{L04243 Gustav Schwarzkopf an Arthur Schnitzler, 16. 6. 1928}
\nopagebreak\mylabel{L04243v}
\rehead{ }\normalsize\beginnumbering\briefempfaengerindex{Schnitzler, Arthur@\textsc{Schnitzler, Arthur}!zzzSchwarzkopf, Gustav@\emph{von Gustav Schwarzkopf}!1928-06-161@{16. 6. 1928}|(be}
\toendnotes[C]{\smallbreak\pagebreak[2]}
\correspDesc{Versand  durch Gustav Schwarzkopf am 16. 6. 1928 in Wien
\newline{}Erhalt  durch Arthur Schnitzler im Zeitraum [17. 6. 1928 – 21. 6. 1928?] in Bad Ischl}\toendnotes[C]{\smallbreak}
\Standort{CUL, Schnitzler, B 96a.}
\physDesc{Bildpostkarte, 503 Zeichen
\newline{}Handschrift: schwarze Tinte, deutsche Kurrent
\newline{}Versand: Stempel: »\nobreak{}\oindex{I., Innere Stadt@\textbf{I., Innere Stadt}, \emph{Verwaltungsgebiet}|pwk}1/\textsubscript{1} Wien
                                       8, 16. VI. 28, 14\nobreak{}«.  }\toendnotes[C]{\smallbreak}\pstart{}{\pb}I. Tiefer Graben 17. II. 6.\oindex{Wien@\textbf{Wien}!I., Innere Stadt@\textbf{I., Innere Stadt}!Tiefer Graben 17@\textbf{Tiefer Graben 17}, \emph{Wohngebäude}|pw}\pend{}{\bigskip}\pstart{}Herrn\pend{}\pstart{}\textsc{D\textsuperscript{r} Arthur Schnitzler}\pend{}\pstart{}\textsc{Bad Ischl\oindex{Bad Ischl@\textbf{Bad Ischl}|pw}}\pend{}\pstart{}\textsc{Hôtel Elisabeth\oindex{Hotel Kaiserin Elisabeth [Bad Ischl]@\textbf{Hotel Kaiserin Elisabeth [Bad Ischl]}, \emph{Hotel}|pw}}\pend{}{\bigskip}\vspace{1em}
\pstart
           \raggedleft{}{\pb}16. VI. 28\pend
           \vspace{0.5em}
\pstart
           lieber Arthur, ich danke Ihnen für Ihre liebe \label{K_L04243-1v}\edtext{Aufforderung}{\lemma{\textnormal{\emph{Aufforderung}}}\Cendnote{\textnormal{nicht überliefert}}}\label{K_L04243-1}. Sie
               haben ja{ }ſo Recht, ich{ }ſollte fortgehen, ich fühle das{ }ſelbſt, aber ich bringe
               vorläufig nicht die Energie auf, zu reiſen. Und da{\geminationn}, Sie
               wollen{ }ſchon am \uline{20} zurückko{\geminationm}en\pend
           
\pstart
           – Heute ist der 16. und das Wetter iſt geſtern und heute{ }ſo{ }ſchlecht.\pend
           
\pstart
           Alſo: Hoffentlich auf Wiederſehen in Wien\oindex{Wien@\textbf{Wien}, \emph{Verwaltungsgebiet}|pw}. Mit der
               Bitte mich Frau P.\pwindex{Pollaczek, Clara Katharina 15.\,1.\,1875 Wien – 22.\,7.\,1951 ebd.@\textsc{Pollaczek, Clara Katharina} (15.\,1.\,1875 Wien – 22.\,7.\,1951 ebd.), \emph{Schriftstellerin}|pw} beſtens zu empfehlen und mit
               herzlichen Grüßen. ihr\pend
           \pstart \spacefill\mbox{Gustav.}\pend{}\selectlanguage{ngerman}\endnumbering\briefempfaengerindex{Schnitzler, Arthur@\textsc{Schnitzler, Arthur}!zzzSchwarzkopf, Gustav@\emph{von Gustav Schwarzkopf}!1928-06-161@{16. 6. 1928}|)be}\mylabel{L04243h}
\begin{anhang}
\end{anhang}\newcommand{\dateiname}{L04243}\newcommand{\titel}{Gustav Schwarzkopf an Arthur Schnitzler, 16. 6. 1928}\newcommand{\editorInnen}{Herausgegeben von Jahnke, SelmaMüller, Martin Anton}\input{../tex-inputs/latex-leseansicht-abspann}
